\documentclass[11pt]{article}
\usepackage[margin=1in]{geometry}
\usepackage{enumitem}
% =============================================================================
% Preambles/header.tex — shared LaTeX/Beamer preamble for this project
%
% Use from a Beamer lecture by prepending:
%     \input{header}
% and compiling with TEXINPUTS=../Preambles:$TEXINPUTS (the /compile-latex
% skill does this automatically).
%
% PALETTE CONTRACT. Color names defined here MUST match the names in
% Quarto/theme-template.scss so Beamer and Quarto renderings use the same
% palette. The scripts/check-palette-sync.sh script enforces this.
%
% To customize: edit the HEX values below AND the matching $variable in
% Quarto/theme-template.scss, then run scripts/check-palette-sync.sh.
% =============================================================================

% -----------------------------------------------------------------------------
% Engine detection + encoding
%
% Our /compile-latex skill uses XeLaTeX, where inputenc is unnecessary and
% can produce encoding warnings. Only load inputenc under pdfTeX.
% -----------------------------------------------------------------------------
\usepackage{iftex}
\ifPDFTeX
  \usepackage[utf8]{inputenc}
\fi

\usepackage{amsmath, amssymb, amsthm}
\usepackage{booktabs}
\usepackage{graphicx}

% -----------------------------------------------------------------------------
% PALETTE — mirrors Quarto/theme-template.scss variable names.
% Load xcolor and define colors BEFORE anything that references them
% (notably hyperref, hypersetup, and the Beamer color assignments).
% -----------------------------------------------------------------------------
\usepackage{xcolor}

\definecolor{primary-blue}{HTML}{012169}      % headings, accents
\definecolor{primary-gold}{HTML}{B9975B}      % emphasis, borders
\definecolor{highlight-yellow}{HTML}{F2A900}  % markers, alerts
\definecolor{light-bg}{HTML}{E8EDF5}          % title / section backgrounds
\definecolor{jet}{HTML}{1A1A1A}               % body text

% Semantic colors (used in diagrams and annotations)
\definecolor{positive}{HTML}{15803D}          % good / observed / identified
\definecolor{negative}{HTML}{B91C1C}          % bad / problematic / bias
\definecolor{neutral}{HTML}{525252}           % reference / context

% Highlight accents (match .hi-* classes in the SCSS)
\definecolor{hi-slate}{HTML}{314F4F}
\definecolor{hi-green}{HTML}{2E8B57}
\definecolor{hi-red}{HTML}{C41E3A}

% Semantic aliases so skill templates and snippets can write
% \textcolor{accent}{...} without hard-coding "primary-blue".
\colorlet{accent}{primary-blue}
\colorlet{accent2}{primary-gold}

% -----------------------------------------------------------------------------
% Hyperref — loaded AFTER xcolor + palette so link colors resolve.
% -----------------------------------------------------------------------------
\usepackage{hyperref}
\hypersetup{colorlinks=true, linkcolor=primary-blue, urlcolor=primary-blue,
            citecolor=primary-blue, pdfborder={0 0 0}}

% -----------------------------------------------------------------------------
% Course code — set once per deck (after \input{header}, before \begin{document})
% via \coursecode{CS401}. Shown in the footer on every slide so a deck's course
% is identifiable without opening the title page. Empty by default.
% -----------------------------------------------------------------------------
\makeatletter
\newcommand{\coursecode}[1]{\gdef\@coursecodetext{#1}}
\gdef\@coursecodetext{}
\makeatother

% -----------------------------------------------------------------------------
% Beamer theme assignments (only apply under Beamer).
%
% We detect Beamer via \@ifclassloaded{beamer}, which is the documented,
% non-internal way. This must be wrapped in \makeatletter/\makeatother
% because \@ifclassloaded contains an @-catcode character.
% -----------------------------------------------------------------------------
\makeatletter
\@ifclassloaded{beamer}{%
  \usefonttheme{professionalfonts}
  \setbeamercolor{structure}{fg=primary-blue}
  \setbeamercolor{frametitle}{fg=primary-blue}
  \setbeamercolor{title}{fg=primary-blue}
  \setbeamercolor{subtitle}{fg=primary-gold}
  \setbeamercolor{author}{fg=primary-blue}
  \setbeamercolor{institute}{fg=neutral}
  \setbeamercolor{date}{fg=primary-gold}
  \setbeamercolor{itemize item}{fg=primary-blue}
  \setbeamercolor{itemize subitem}{fg=primary-gold}
  \setbeamercolor{alerted text}{fg=primary-gold}
  \setbeamercolor{block title}{fg=white, bg=primary-blue}
  \setbeamercolor{block body}{bg=light-bg}
  % Semantic block variants: block=definition/concept (blue), exampleblock=
  % worked example (green/positive), alertblock=key takeaway or caution (gold).
  % Keep to <=2 colored boxes per slide across ALL three variants (INV-7).
  \setbeamercolor{block title example}{fg=white, bg=positive}
  \setbeamercolor{block body example}{bg=light-bg}
  \setbeamercolor{block title alerted}{fg=white, bg=primary-gold}
  \setbeamercolor{block body alerted}{bg=light-bg}

  % Minimal footer: course code (left) + page X / N (right), no navigation clutter
  \setbeamertemplate{navigation symbols}{}
  \setbeamertemplate{footline}{%
    \color{neutral}\footnotesize\hspace{0.7em}\@coursecodetext\hfill
    \insertframenumber/\inserttotalframenumber\hspace{0.7em}\vspace{0.3em}}
}{}
\makeatother

% -----------------------------------------------------------------------------
% TikZ setup — libraries the snippet gallery uses
% -----------------------------------------------------------------------------
\usepackage{tikz}
\usetikzlibrary{arrows.meta, positioning, calc, decorations.pathreplacing,
                fit, shapes.geometric, backgrounds}

% Shared TikZ styles — snippets and hand-written diagrams can pick these up
% via `\begin{tikzpicture}[every node/.style={...}]` or by naming specific
% styles (e.g., `\node[dag-node] (X) at (0,0) {X};`).
\tikzset{
  % Node styles for causal DAGs and similar
  dag-node/.style = {
    draw=primary-blue, fill=light-bg, rounded corners=2pt,
    minimum width=1.2cm, minimum height=0.9cm, align=center, font=\small
  },
  decision-node/.style = {
    draw=primary-gold, fill=white, diamond, aspect=1.8,
    minimum width=1.8cm, minimum height=1.2cm, align=center, font=\small,
    inner sep=1pt
  },
  % Edge styles — observed vs counterfactual
  observed-edge/.style = {-{Latex[length=2.5mm]}, thick, draw=jet},
  counterfactual-edge/.style = {-{Latex[length=2.5mm]}, thick, draw=neutral, dashed},
  confound-edge/.style = {-{Latex[length=2.5mm]}, thick, draw=negative, dashed},
  % Data-point styles for plots
  observed-dot/.style = {fill=primary-blue, draw=primary-blue, circle, inner sep=1.2pt},
  counterfactual-dot/.style = {fill=white, draw=neutral, circle, inner sep=1.2pt}
}

% -----------------------------------------------------------------------------
% Convenience macros
% -----------------------------------------------------------------------------
\newcommand{\muted}[1]{\textcolor{neutral}{#1}}
\newcommand{\key}[1]{\textcolor{primary-gold}{\textbf{#1}}}
\newcommand{\good}[1]{\textcolor{positive}{#1}}
\newcommand{\bad}[1]{\textcolor{negative}{#1}}

% Standout frame macro: full-bleed dark background for section transitions.
% Beamer-only (uses \begin{frame}/\setbeamercolor/\usebeamerfont) -- do not
% call from an article-class file (e.g. Notes/<CODE>/*-notes.tex). \muted,
% \key, \good, \bad, and \coursecode above are plain xcolor/gdef and are
% safe to use from either class; verified when Notes/article-class support
% was added.
% Expand: \transitionslide{Your transition title}
\newcommand{\transitionslide}[1]{%
  \begingroup
  \setbeamercolor{background canvas}{bg=primary-blue}
  \begin{frame}[plain]
    \centering
    \vfill
    {\usebeamerfont{title}\color{white}\Large #1\par}
    \vfill
  \end{frame}
  \endgroup
}

% Section divider frame (Beamer-only), an alternative to \transitionslide with a
% darker two-line style: near-black (jet) background, a small white label line
% above a larger gold title, and a thin gold rule along the bottom edge.
% Expand: \sectiondivider{Section 2}{Pointers}
\newcommand{\sectiondivider}[2]{%
  \begingroup
  \setbeamercolor{background canvas}{bg=jet}
  \begin{frame}[plain]
    \vspace*{3.0cm}
    {\color{white}\ttfamily\large #1\par}
    \vspace{0.5cm}
    {\color{primary-gold}\LARGE #2\par}
    \begin{tikzpicture}[remember picture, overlay]
      \draw[primary-gold, line width=2.5pt]
        ($(current page.south west)+(0,0.1)$) -- ($(current page.south east)+(0,0.1)$);
    \end{tikzpicture}
  \end{frame}
  \endgroup
}

% -----------------------------------------------------------------------------
% Channel watermark — "Concepts That Click" (YouTube).
%
% Article-class files (CompetitiveExam Q/A sets, Notes, Assignments) get a
% light diagonal draft watermark on every page plus a centered footer naming
% the channel. Beamer decks get a subtle TikZ background watermark instead
% (the footline already carries the course code). Centralized here so every
% MCQ PDF is branded without per-file edits.
% -----------------------------------------------------------------------------
\makeatletter
\@ifclassloaded{article}{%
  \usepackage{draftwatermark}
  \SetWatermarkText{Concepts That Click}
  \SetWatermarkScale{1}
  \SetWatermarkFontSize{2cm}
  \SetWatermarkLightness{0.86}
  \SetWatermarkAngle{45}
  \usepackage{fancyhdr}
  \pagestyle{fancy}
  \fancyhf{}
  \fancyfoot[C]{\footnotesize\textcolor{neutral}{Concepts That Click~|~YouTube: \href{https://www.youtube.com/@concepts.thatclick}{@concepts.thatclick}}}
  \renewcommand{\headrulewidth}{0pt}
  \renewcommand{\footrulewidth}{0pt}
  \fancypagestyle{plain}{%
    \fancyhf{}%
    \fancyfoot[C]{\footnotesize\textcolor{neutral}{Concepts That Click~|~YouTube: \href{https://www.youtube.com/@concepts.thatclick}{@concepts.thatclick}}}%
    \renewcommand{\headrulewidth}{0pt}%
    \renewcommand{\footrulewidth}{0pt}%
  }
}{}
\@ifclassloaded{beamer}{%
  \setbeamertemplate{background}{%
    \begin{tikzpicture}[remember picture, overlay]
      \node[rotate=45, opacity=0.055, align=center]
        at (current page.center)
        {\Huge\textcolor{neutral}{Concepts That Click}};
    \end{tikzpicture}%
  }
}{}
\makeatother 
\coursecode{JKSSB}
\renewcommand{\thesection}{42.\arabic{section}}
\newtheorem{example}{Example}[section]
\newenvironment{definitionbox}[1]{%
  \par\noindent\textbf{Definition (#1).}\ \itshape
}{\par\vspace{0.3em}}
\title{Malware --- Detailed Lecture Notes}
\author{JKSSB Computer Awareness}
\date{\today}

\begin{document}
\maketitle

\section{What ``malware'' means}
Malware is short for \textbf{malicious software}: any program written to
harm a computer, steal from it, or misuse it. The official syllabus names
\textbf{computer virus and latest anti-virus} explicitly, and the exam's
computer block tests the near neighbours as well: worm, Trojan, ransomware,
spyware, adware, rootkit, keylogger, bot and botnet, backdoor, logic bomb,
and cryptojacking.

\begin{definitionbox}{Malware}
Any software designed to damage a computer, steal data from it, or misuse
its resources without the owner's consent.
\end{definitionbox}

\begin{center}
\begin{tabular}{p{0.24\textwidth}p{0.66\textwidth}}
\toprule
\textbf{Term} & \textbf{Meaning in this lesson} \\
\midrule
Virus & Malware that attaches to a host file and needs human action to run. \\
Worm & Standalone, self-spreading malware that travels over networks. \\
Trojan (Trojan horse) & Malware disguised as useful software; does not replicate itself. \\
Ransomware & Malware that locks files and demands payment for release. \\
Spyware & Malware that secretly watches the user and steals data. \\
Adware & Malware that floods the machine with unwanted advertisements. \\
Keylogger & A spying tool that records every key pressed. \\
Rootkit & Malware that hides itself and other malware deep in the system. \\
Backdoor & A secret entry path into the system bypassing normal checks. \\
Bot / botnet & A remotely controlled machine / a network of such machines. \\
Logic bomb & Malicious code that waits for a trigger, then strikes. \\
Cryptojacking & Secret use of the victim's computer to mine cryptocurrency. \\
Anti-virus & Software that detects, blocks, and removes known malware. \\
\bottomrule
\end{tabular}
\end{center}

The rest of this lesson uses these terms in one consistent sense. A
\textbf{virus} always needs a host file and a human launch. A \textbf{worm}
needs neither. A \textbf{Trojan} pretends to be legitimate and does not
replicate. \textbf{Ransomware} encrypts and demands money.
\textbf{Spyware} hides and steals; \textbf{adware} shows ads. A
\textbf{keylogger} records keystrokes. A \textbf{rootkit} conceals; a
\textbf{backdoor} re-admits. A \textbf{botnet} is a controlled network (robot
+ network). A \textbf{logic bomb} waits for its trigger.
\textbf{Cryptojacking} mines on the victim's machine.

\section{Virus, worm, and Trojan}
A \textbf{virus} attaches itself to a host program or file and stays inert
until a person opens, runs, or shares that file. Common carriers are email
attachments, pen drives, and downloaded files. Once launched it can damage,
delete, or corrupt files. The defining limit is dependence: no host file
and no human action, no spread.

\begin{definitionbox}{Virus}
Malware that attaches to a host program or file and requires human action
to execute and spread.
\end{definitionbox}

A \textbf{worm} is standalone and self-replicating. It copies itself from
machine to machine over network connections, exploiting security gaps,
without any host file and without any user click. Its main harm is
consumption: networks clog, systems slow, and services jam under the load
of endless self-copies.

\begin{definitionbox}{Worm}
Standalone malware that replicates and spreads across networks by itself,
without a host file or human action.
\end{definitionbox}

A \textbf{Trojan}, or Trojan horse, takes its name from deception: it
arrives looking like something useful --- a free game, a handy tool, a
harmless document. The user installs or opens it willingly, and the hidden
malicious part runs alongside. Unlike a virus or a worm, a Trojan does not
replicate itself; each new machine needs its own act of trickery. Once
inside, it can steal data, install spyware or ransomware, or open a backdoor
for the attacker's later return.

\begin{definitionbox}{Trojan (Trojan horse)}
Malware disguised as legitimate software that tricks the user into
installing it; it does not replicate itself.
\end{definitionbox}

The three are best remembered as one table: a virus needs a host file and a
user action; a worm needs neither; a Trojan needs the user's trust and does
not self-replicate. File damage points to a virus, network-wide slowdown
with no clicks points to a worm, and a ``useful'' download that betrays the
machine points to a Trojan.

\begin{example}
An item describes a program that ``spreads from computer to computer over
the network without any user opening a file.'' Trace the table: needs no
host and no click means worm. If instead the stem says the malware ``was
installed as part of a free utility the user downloaded,'' the deception
plus no replication means Trojan.
\end{example}

\section{Ransomware}
\textbf{Ransomware} locks the victim out of their own data, usually by
encrypting files (or the whole screen), and then displays a demand: pay the
ransom to receive the unlock key. The ``pay to unlock'' message, often with
a countdown or a payment address, is its signature.

Ransomware usually arrives the Trojan way --- a harmful attachment or link
the user opens --- and then encrypts documents, photos, and databases in
minutes. The expected defences are regular \textbf{offline backups}, updated
operating system and \textbf{anti-virus} software, and
caution with attachments. The exam's standing line is that paying the ransom
does not guarantee recovery; the backup is what defeats the demand.

\begin{definitionbox}{Ransomware}
Malware that encrypts the victim's files (or locks the screen) and demands
payment for the key to restore them.
\end{definitionbox}

\section{Spyware, adware, and the keylogger}
\textbf{Spyware} is the silent watcher. It hides on the machine, observes
browsing habits, logins, and personal details, and sends them to a third
party. It deliberately causes no visible damage so the victim suspects
nothing. It commonly arrives bundled with free downloads or deceptive
installers, and its purpose is theft by observation: passwords, bank
details, and identity data.

\textbf{Adware} is the noisy opposite. It floods the screen with unwanted
advertisements --- pop-ups, banners, redirected searches --- to earn money
from ad traffic. Heavy adware also slows the machine and may carry spyware
along with it, so the exam treats persistent pop-ups as more than a
nuisance.

A \textbf{keylogger} is a focused spying tool: it records every key pressed
and forwards the log to the attacker, harvesting usernames, passwords, and
card numbers. Keyloggers exist as software and as small hardware devices
fitted between the keyboard and the computer.

\begin{definitionbox}{Spyware vs adware}
Spyware hides and steals information quietly; adware displays unwanted
advertisements openly. A keylogger is a spying tool that records keystrokes.
\end{definitionbox}

\begin{example}
An item reports ``constant pop-up advertisements and redirected searches.''
The visible ads mean adware. If instead the item reports ``passwords stolen
although nothing on screen looked wrong,'' the silent theft means spyware,
and if it specifies ``every keystroke was recorded,'' the tool is a
keylogger.
\end{example}

\section{Rootkit and backdoor}
A \textbf{rootkit} is malware whose job is concealment. It lodges deep in
the system, often with administrator-level control, and hides files,
processes, and network activity from the user and from ordinary detection
tools --- including hiding other malware.

A \textbf{backdoor} is the secret entry a rootkit or Trojan protects: a
path into the system that bypasses normal login and security checks, left
open so the attacker can return at will. The pairing to remember is that
Trojans and rootkits commonly install backdoors; the rootkit hides the
intruder, and the backdoor keeps the door open.

\begin{definitionbox}{Rootkit and backdoor}
A rootkit conceals an intruder and its malware deep in the system; a
backdoor is a hidden entry path that bypasses normal authentication.
\end{definitionbox}

\section{Bot and botnet, logic bomb, cryptojacking}
A \textbf{bot} is an ordinary computer that has been compromised and is now
obeyed remotely; a \textbf{botnet} is a network of such machines (the word
joins robot + network) under one master. Botnets send spam, spread further
malware, and launch large-scale attacks that no single machine could mount.
The owner typically notices only slowness.

A \textbf{logic bomb} is malicious code already planted that waits for a
trigger --- a date, an event, or a specific action --- and strikes only when
the condition is met. It does not spread; it waits.

\textbf{Cryptojacking} is theft of computing power: the attacker secretly
runs cryptocurrency mining on the victim's machine, so the computer runs hot
and slow while the attacker collects the earnings. Like a bot, its outward
sign is an otherwise unexplained slowdown.

\section{How malware reaches a machine, and how to stop it}
Malware reaches users through a short list of routes: harmful email
attachments and links, infected pen drives, bundled free-software
installers, downloads from untrusted sites, and unpatched system software
with known gaps.

The defences form a matching list. \textbf{Anti-virus} software detects,
blocks, quarantines, and removes known malware; both it and the operating
system must be kept updated, since each update closes gaps and refreshes the
list of known threats. A \textbf{firewall} filters incoming and outgoing
network traffic and blocks suspicious connections, but it is a barrier, not
a cleaner: it does not remove malware already inside. Safe habits complete
the set: do not open unknown attachments or links, scan pen drives before
use, download only from trusted sources, and keep regular \textbf{offline
backups} so ransomware loses its leverage.

\section{Exam retrieval and common confusions}
Six distinctions prevent most errors on this topic.
\begin{enumerate}
  \item A virus needs a host file \textbf{and} a human action; a worm needs
    \textbf{neither}; a Trojan disguises itself and \textbf{does not
    replicate}.
  \item Spyware \textbf{hides and steals}; adware \textbf{shows ads}. A
    keylogger records \textbf{keystrokes}.
  \item Ransomware \textbf{encrypts and demands payment}; its signature is
    the ``pay to unlock'' message, and the defence is the \textbf{offline
    backup}.
  \item A rootkit \textbf{hides} the intruder; a backdoor is the
    \textbf{secret re-entry}. Trojans commonly install backdoors.
  \item A botnet is a \textbf{network} of remotely controlled machines
    (robot + network), used for spam and large attacks.
  \item A logic bomb \textbf{waits for a trigger}; cryptojacking
    \textbf{mines cryptocurrency} on the victim's machine.
  \item Anti-virus \textbf{removes} malware; a firewall \textbf{filters}
    traffic and does not clean an infected machine.
\end{enumerate}

\begin{example}
An item asks which malware ``encrypts the user's files and demands money
for the decryption key.'' The lock plus the demand means ransomware ---
not a worm (self-spreading), not adware (ads), not spyware (silent theft).
If the stem instead says the program ``waits until a fixed date and then
deletes files,'' the waiting trigger means a logic bomb.
\end{example}

\subsection*{Exercises}
\begin{enumerate}[label=\arabic*.]
  \item Distinguish a virus from a worm from a Trojan on three points: host
    file, self-spreading, and user action.
  \item Contrast spyware with adware. Where does a keylogger fit?
  \item Explain what ransomware does and name the defence that defeats its
    demand.
  \item Distinguish a rootkit from a backdoor, and explain how they work
    together.
  \item Define bot/botnet, logic bomb, and cryptojacking in one sentence
    each, naming each type's outward sign.
\end{enumerate}

\subsection*{Solutions}
\begingroup\small
\begin{enumerate}[label=\arabic*.]
  \item A virus needs a host file and a human launch and does not spread by
    itself; a worm is standalone and self-spreading over networks with no
    host and no click; a Trojan pretends to be legitimate software, is
    installed willingly, and does not replicate itself.
  \item Spyware hides and secretly steals user data; adware openly floods
    the machine with unwanted ads. A keylogger is a focused spying tool that
    records every keystroke to steal credentials.
  \item Ransomware encrypts the victim's files (or locks the screen) and
    demands payment for the unlock key. Regular offline backups defeat the
    demand, since the data can be restored without paying.
  \item A rootkit hides the intruder and its malware deep in the system; a
    backdoor is a secret entry bypassing normal checks. Trojans and rootkits
    install backdoors so the attacker can return unseen.
  \item A bot is a remotely controlled compromised machine and a botnet is a
    network of them (sign: slowness during spam or attacks); a logic bomb
    waits for a trigger such as a date and then strikes (sign: a timed
    surprise); cryptojacking secretly mines cryptocurrency on the victim's
    machine (sign: hot, slow computer).
\end{enumerate}
\endgroup
\end{document}
